\thispagestyle{plain}
\begin{center}
{\Large \bf Abstract}
\end{center}
\vspace{0.6cm}

\begin{justify}
This report documents a comprehensive project carried out during a full-time engineering internship at Ignosis, a fintech infrastructure company in Ahmedabad. The brief was to design and contribute to an AI-Based Collections Platform that is being used by banks, non-banking financial companies and digital-lending fintechs in India for the recovery of overdue retail loan accounts.

In its present form the collections process for an Indian retail lender is largely manual. Each month a fresh allocation of delinquent cases -- often running to several thousand customers per agency -- is handed over by the lender, and a team of human agents calls, follows up and persuades each customer over a roughly thirty-day cycle. The agency is paid only on recovery, so the entire business model is sensitive to contact rate, agent productivity and the speed with which the right customer is reached on the right channel. Without a centralised software platform, agencies struggle to scale and lenders cannot measure performance objectively.

The platform that this project contributed to is a single web-based application built around three interlocking ideas. At the centre is a tenant-configurable {\bf workflow engine} -- a small expression-driven graph of nodes and transitions -- that captures the lender's collection policy as data rather than code. A {\bf campaign module} wraps thousands of customer cases into a single executable unit that is driven through the workflow. An {\bf AI voice calling engine}, built on the firm's internal voice AI platform and a telephony gateway, becomes one of the workflow's action nodes, which lets the platform reach customers outside human agent working hours without any change to the rest of the system. A separate {\bf multi-channel automation layer} sends WhatsApp messages from real consumer-grade SIMs and dispatches SMS from a custom Android application, again as workflow actions.

The contributions described in this report span the case detail and case listing pages of the agent dashboard, the end-to-end campaign creation and CSV upload pipeline, the workflow nodes that integrate AI calling into a campaign, the speed dialer that auto-loads the next case after each disposition, the WhatsApp service with multi-device rotation and proactive block-avoidance, the desktop wrapper that exposes these capabilities to non-technical agents, and the Android SMS sender app. The platform is reaching a Bucket~0 collection rate of about eighty-five percent in production and is recovering close to seventy percent of allocated cases within the first five days of each monthly cycle, both of which are material improvements over the agency-only baseline.
\end{justify}

\clearpage
